% !TEX root = ../../main.tex
%
% All-tier Virasoro residues: construction and comparison obligations.

\providecommand{\Vir}{\mathrm{Vir}}
\providecommand{\cA}{\mathcal A}
\providecommand{\ClaimStatusDefinitional}{}
\providecommand{\ClaimStatusConditional}{}
\providecommand{\ClaimStatusOpen}{}

\chapter{All-Tier Virasoro Residues}
\label{ch:all-tier-generating-function}

\section*{The object to be constructed}

The Virasoro operator product and the level-four vacuum Gram form
determine exact local data.  In the notation of
Chapter~\ref{ch:shadow-tower-higher-coefficients},
\[
 T(z)T(w)\sim \frac{c/2}{(z-w)^4}
 +\frac{2T(w)}{(z-w)^2}+\frac{\partial T(w)}{z-w},
 \qquad
 N_\Lambda(c)=\frac{c(5c+22)}{10}.
\]
The connected Ward functions are rational functions on ordered
configuration spaces.  Scalar higher residues arise after the
ordered-residue package
\(\mathsf H_{\mathrm{res}}(\Vir_c;X)\) supplies a completed complex,
a collision-compatible correlator map, and contractions onto the
scalar lines.

The all-tier problem begins with this chain model.  A formal
weighted-Riccati series is an auxiliary algebraic object.  Its
identification with the scalar coordinates
\(S_r(\Vir_c;\mathsf H_{\mathrm{res}})\) is a comparison theorem to be
proved from the transfer data.

\section{Laurent data of a constructed residue tower}

\begin{definition}[Asymptotic package and normalized tiers]
\ClaimStatusDefinitional
\label{def:bivariate-F}
Assume \(\mathsf H_{\mathrm{res}}(\Vir_c;X)\).  The package
\(\mathsf H_{\mathrm{asymp}}\) consists of:
\begin{enumerate}[(i)]
\item meromorphic dependence of every transferred coordinate on the
central charge~\(c\);
\item a declared normalization exponent \(\nu_r\) and a Laurent
expansion at \(c=\infty\);
\item uniform coefficient bounds in~\(r\), together with a declared
order of the limits \(c\to\infty\) and \(r\to\infty\);
\item compatibility of coefficient extraction with the filtration
completion used in \(\mathsf H_{\mathrm{res}}\).
\end{enumerate}
Write
\begin{equation}
 \Phi_r(c;\mathsf H_{\mathrm{res}})
 :=c^{\nu_r}S_r(\Vir_c;\mathsf H_{\mathrm{res}})
 =\sum_{K\ge1}T_K(r)c^{1-K}.
 \label{eq:all-tier-raw-intro}
\end{equation}
The leading tier is \(A_r:=T_1(r)\).  Whenever \(A_r\) is invertible,
set
\begin{equation}
 F(r,u):=\frac{\Phi_r(1/u;\mathsf H_{\mathrm{res}})}{A_r}
 =1+\sum_{K\ge2}\frac{T_K(r)}{A_r}u^{K-1}.
 \label{eq:bivariate-F-intro}
 \label{eq:F-def}
\end{equation}
Thus the bivariate series is an invariant of
\((\Vir_c,\mathsf H_{\mathrm{res}},\mathsf H_{\mathrm{asymp}})\).
\end{definition}

\begin{remark}[Transferred equation before scalar closure]
Write \(s=s_{\Vir_c}\) for the transferred scalar Maurer--Cartan
element.  Its coordinates satisfy the projected transferred
Maurer--Cartan equation
\begin{equation}
 \operatorname{pr}_r\left(
 \sum_{n\ge1}\frac1{n!}
 \ell_n^{\mathrm{sc}}(s,\ldots,s)\right)=0.
 \label{eq:bilinear-master-intro}
\end{equation}
A bilinear recurrence follows after a quadratic-closure package
identifies the surviving transferred brackets and their kernels.
This package determines the admissible arity pairs and every scalar
normalization.
\end{remark}

\section{Generating-function and tier problems}

\begin{problem}[Construct the bivariate residue generating function]
\ClaimStatusOpen
\label{thm:all-tier-bivariate-generating-function}
\label{thm:all-tier-laurent-stratification}
\label{thm:all-tier-closed-form-proved}
\emph{Type signature:}
\textup{(Open quadrant, completed ordered residue-bar and transferred
scalar presentations, level \(2\),
\(\mathsf H_{\mathrm{res}}+\mathsf H_{\mathrm{quad}}
+\mathsf H_{\mathrm{asymp}}\)).}

Construct the Virasoro residue package and determine the transferred
brackets in~\eqref{eq:bilinear-master-intro}.  Then derive the
functional equation of \(F(r,u)\) directly from those brackets.
A complete solution supplies:
\begin{enumerate}[(i)]
\item chain representatives for every scalar coordinate;
\item the quadratic and higher transfer kernels;
\item a proof of the Laurent expansion in
Definition~\ref{def:bivariate-F};
\item a closed form, when one exists, derived from the constructed
functional equation;
\item an independent residue computation of the same coefficients.
\end{enumerate}
Any hypergeometric expression belongs to the conclusion of this
comparison theorem.
\end{problem}

\begin{problem}[Determine the tier arithmetic]
\ClaimStatusOpen
\label{cor:tier-k-leading-coefficient}
\label{thm:tier-5-closed-form}
\label{cor:tier-leading-denominator-pattern}
\label{cor:tier-K-kummer-arithmetic}
\emph{Type signature:}
\textup{(Open quadrant, scalar Laurent presentation, level \(2\),
\(\mathsf H_{\mathrm{res}}+\mathsf H_{\mathrm{asymp}}
+\mathsf H_{\mathrm{arith}}\)).}

For the constructed coefficients \(T_K(r)\), determine the leading
polynomial degree, reduced denominator, and prime support.  The
arithmetic package \(\mathsf H_{\mathrm{arith}}\) records exact
integer normalizations, factorization certificates, and a second
coefficient oracle.  Kummer comparisons are then intersections of
explicitly certified prime sets.
\end{problem}

\begin{problem}[Derive the analytic differential equation]
\ClaimStatusOpen
\label{thm:all-tier-fuchsian-ode}
\label{cor:three-disjoint-hz-iv-chains}
\emph{Type signature:}
\textup{(Open quadrant, analytic scalar presentation, level \(2\),
\(\mathsf H_{\mathrm{res}}+\mathsf H_{\mathrm{asymp}}
+\mathsf H_{\mathrm{ODE}}\)).}

Determine the singular locus and minimal differential operator of
the constructed function \(F(r,u)\).  The package
\(\mathsf H_{\mathrm{ODE}}\) consists of a convergence domain, an
operator derivation from the residue functional equation, and a
Frobenius comparison at every singular point.  A Gauss or Heun
description becomes a theorem after this derivation.
\end{problem}

\section{Comparison with formal algebraic lanes}

\begin{remark}[Formal series and geometric residues]
\label{rem:F-reduction-structure}
A square-root or Catalan series defined from chosen seeds has exact
formal coefficients.  A comparison morphism from its coefficient
complex to the ordered-residue scalar model determines the locus on
which those coefficients compute
\(S_r(\Vir_c;\mathsf H_{\mathrm{res}})\).
The comparison includes the arity convention, residue orientation,
and gauge compatibility.
\end{remark}

\begin{remark}[Finite-tier recovery]
\label{rem:four-tier-recovery}
The first residue tiers are the worked cases for the all-tier
problem.  Their computation begins with the exact Ward correlators
and ends with the normalized contractions in
\(\mathsf H_{\mathrm{res}}\).
\end{remark}

\begin{remark}[Arithmetic signatures]
\label{rem:tier-5-numerical-signature}
\label{rem:fuchsian-heun-comparison}
Certified numerical signatures belong to
\(\mathsf H_{\mathrm{arith}}\); analytic singularities belong to
\(\mathsf H_{\mathrm{ODE}}\).  Agreement between these packages is a
cross-check on the common residue tower.
\end{remark}

\begin{remark}[Analytic radius and adjacent families]
\label{rem:borel-analytic-continuation}
\label{rem:hypergeometric-contiguous-relation}
\label{rem:bivariate-generalisation-W-N}
The analytic radius follows from the singularities of the constructed
generating function.  Principal \(\mathcal W_N\) families require
multivariable residue complexes for their additional generator lines;
their scalar diagonals and mixed channels carry separate transfer
kernels.
\end{remark}

%% Legacy equation anchors are retained for downstream references.  Their
%% mathematical content is governed by the open packages above.
\phantomsection\label{eq:F-hyper}%
\phantomsection\label{eq:master-bilinear-proof}%
\phantomsection\label{eq:master-in-u}%
\phantomsection\label{eq:F-clear-denominator}%
\phantomsection\label{eq:tier-K-closed-form}%
\phantomsection\label{eq:F-sum-over-m}%
\phantomsection\label{eq:tier-5-explicit}%
\phantomsection\label{eq:fuchsian-ode}%
\phantomsection\label{eq:fuchsian-ode-expanded}%
\phantomsection\label{eq:fuchsian-sing-loci}%
\phantomsection\label{eq:gauss-ode-x}%
